\clearpage
\section{问题三参数、补测概率与复现材料}\label{app:q3repro}
\subsection{补测的中心射线近似}\label{app:q3prob}
设首个有信号测站为 $S_0$，示向度为 $\theta_0$。沿 $x(t)=S_0+t\,u(\theta_0)$ 在 $5\sim1500$ m 上等距取 151 个样本。用极坐标面积元中的 $t$ 为位置权重，并采用接收半径在 $[1000,1500]$ 上均匀的建模近似。正观测要求 $\rho\ge t$，历史负观测 $S_j$ 要求 $\rho<\|x(t)-S_j\|$，因此
\[
 \ell(t)=\max(1000,t),\qquad
 H(t)=\min\{1500,\min_{j\in I^-}\|x(t)-S_j\|\}.
\]
$I^-$ 是负观测索引；为空时内层最小值取正无穷。令 $[a]_+=\max(a,0)$，分母把可能的接收半径区间长度求和，分子进一步要求能在候选站 $p$ 接收：
\begin{equation}\label{eq:q3-prob}
 \widehat P_{\rm recv}(p)=
 \frac{\sum_t t\,\mathbf1_\Omega(x(t))
 [H(t)-\max\{\ell(t),\|x(t)-p\|\}]_+}
 {\sum_t t\,\mathbf1_\Omega(x(t))[H(t)-\ell(t)]_+}.
\end{equation}
分母为零时不施加此概率门槛，交由后续几何筛选。它是中心射线上的离散近似，不是完整角楔上的贝叶斯后验。单方位点估计另取 301 个射线样本，剔除靶区外及历史负观测 1000 m 圆盘内样本，作 $t$ 加权平均；该操作不用于缩小保守外包。

\subsection{算法参数与实验设置}
\begin{table}[H]\centering
\caption{问题三算法参数}
\begin{tabular}{@{}p{.32\textwidth}p{.60\textwidth}@{}}\toprule
参数 & 默认取值与含义\\\midrule
角楔半宽、圆外包 & $(1.005+10^{-9})^\circ$，每圆 32 条切线\\
清除及近场半径 & 20 m、5 m；近场仍等待成功清除回复\\
field 距离系数 & $\lambda=1.5$，100 m 行动候选网格\\
sweep / tour 主圈 & $(1100\ \mathrm{m},8)$ / $(1000\ \mathrm{m},8)$\\
adaptive 公共基线 & 400 m，72 个候选方向\\
沿途逼近 & 接近比例 0.55，侧移上限 100 m，最多 12 步\\
连续证书 & $h=50$ m，$r_{\rm cert}=964.644660\ldots$ m\\
joint 覆盖候选 & 1050、1200、1350、1500 m 四圈，各 72 方向\\
贪心覆盖 / 规划候选 & 最多 30 次；$(\kappa,\nu)$ 取 $(1,0.3)$、$(1,1)$、$(2,0.3)$、$(2,1)$、$(3,0.3)$\\
测站精修 & 删除冗余站、保留独占格移动；最多 3 轮\\
中心试探 & 至少两条方位，$r_c\le400$ m；失败后原地补测\\
可选补测 & 概率门槛 0.5、交会基线锐角至少 $25^\circ$\\
清除后扫描 & \texttt{after\_service\_scan='duty'}，实际终点更新\\
小规模精确排序 & 不超过 3 个路点枚举；其余 NN+2-opt\\
离线虚拟时钟上限 & 360000 s；程序墙钟耗时另记，非虚拟成本\\\bottomrule
\end{tabular}\end{table}
各策略使用自身默认配置；最终两代的有效覆盖半径为证书半径。构造性自检不参与性能统计。

\newpage
\subsection{补充统计}
随机场景的源数均匀取 $10\sim16$，源位按圆域面积均匀分布；外圈和中心组均固定 12 源，其半径分别在 $1300\sim1800$ m 和 $0\sim260$ m 内均匀取值。接收半径通常在 $[1000,1500]$ m 内均匀分布，最小接收组固定为 1000 m；非平滑组采用有界位置哈希误差，同地误差固定且不超过 $1^\circ$。按正文场景顺序编号 $s=0,\ldots,4$，第 $i$ 组种子为 $2026091200+100000s+97i$，$i$ 从 0 起计。

下表时间单位为 s，行程为 km；每源列为逐组比值平均，计算列为策略运行的本地实测耗时，不含场景生成与数据分析。
\begingroup
\setlength{\tabcolsep}{3pt}
\begin{longtable}{@{}lrrrrrrrr@{}}
\caption{全部场景的完整成本与清除结果}\\
\toprule 场景/策略 & 清除比例 & 全清 & 时间均值 & P90 & 每源 & 行程 & 检测 & 计算\\\midrule\endfirsthead
\toprule 场景/策略 & 清除比例 & 全清 & 时间均值 & P90 & 每源 & 行程 & 检测 & 计算\\\midrule\endhead
random/field & 0.995 & 47/50 & 5347.3 & 5820.3 & 422.2 & 21.04 & 180.5 & 0.040 \\
random/sweep & 0.986 & 43/50 & 4186.4 & 4609.7 & 333.9 & 16.78 & 128.1 & 0.046 \\
random/tour & 1.000 & 50/50 & 4028.2 & 4673.7 & 316.9 & 14.41 & 177.1 & 0.148 \\
random/adaptive & 1.000 & 50/50 & 3347.9 & 3601.6 & 263.6 & 11.72 & 156.7 & 0.228 \\
random/joint & 1.000 & 50/50 & 3066.7 & 3366.7 & 241.1 & 11.17 & 127.4 & 0.707 \\
annulus/field & 0.996 & 19/20 & 5015.5 & 5306.7 & 419.7 & 18.68 & 204.3 & 0.038 \\
annulus/sweep & 0.975 & 15/20 & 4397.8 & 5029.2 & 377.8 & 16.93 & 159.2 & 0.041 \\
annulus/tour & 1.000 & 20/20 & 4376.3 & 4774.1 & 364.7 & 15.19 & 210.3 & 0.125 \\
annulus/adaptive & 1.000 & 20/20 & 3598.7 & 3945.9 & 299.9 & 12.52 & 172.8 & 0.195 \\
annulus/joint & 1.000 & 20/20 & 3215.8 & 3455.7 & 268.0 & 11.43 & 145.6 & 0.472 \\
center/field & 1.000 & 10/10 & 4293.7 & 4379.0 & 357.8 & 16.05 & 170.9 & 0.040 \\
center/sweep & 1.000 & 10/10 & 2632.3 & 2772.4 & 219.4 & 9.37 & 116.3 & 0.055 \\
center/tour & 1.000 & 10/10 & 2529.6 & 2759.1 & 210.8 & 9.23 & 100.3 & 0.110 \\
center/adaptive & 1.000 & 10/10 & 2573.5 & 2675.7 & 214.5 & 8.23 & 144.8 & 0.179 \\
center/joint & 1.000 & 10/10 & 2256.1 & 2377.0 & 188.0 & 8.01 & 98.8 & 0.885 \\
hash/field & 0.992 & 18/20 & 5451.7 & 5993.3 & 425.3 & 21.53 & 180.3 & 0.041 \\
hash/sweep & 0.977 & 17/20 & 4060.3 & 4552.1 & 324.3 & 16.24 & 125.0 & 0.050 \\
hash/tour & 1.000 & 20/20 & 4012.8 & 4730.6 & 308.7 & 14.73 & 163.1 & 0.157 \\
hash/adaptive & 1.000 & 20/20 & 3303.9 & 3587.9 & 256.2 & 11.67 & 150.7 & 0.236 \\
hash/joint & 1.000 & 20/20 & 3047.0 & 3261.9 & 236.6 & 11.22 & 121.5 & 0.734 \\
worstrecv/field & 1.000 & 20/20 & 4854.4 & 5105.9 & 381.0 & 18.22 & 191.9 & 0.035 \\
worstrecv/sweep & 0.861 & 7/20 & 4465.0 & 5011.3 & 410.8 & 17.16 & 163.7 & 0.038 \\
worstrecv/tour & 1.000 & 20/20 & 4438.8 & 4987.5 & 346.5 & 15.66 & 204.5 & 0.163 \\
worstrecv/adaptive & 1.000 & 20/20 & 3582.2 & 3811.7 & 280.9 & 12.14 & 182.1 & 0.215 \\
worstrecv/joint & 1.000 & 20/20 & 3320.7 & 3565.5 & 260.4 & 11.79 & 149.9 & 0.630 
\\
\bottomrule\end{longtable}
\endgroup
\newpage
四项消融只用随机集前 20 局，完整策略均值 \QThreeAblationBase\ s。下表增长率定义为“消融均值/完整策略均值$-1$”，区间是同一组种子重采样的 4000 次配对 bootstrap，未进行多重比较校正；结论限于这组机制比较。
\begin{table}[H]\centering
\caption{单因素消融的成本变化及配对区间}
\begin{tabular}{@{}lrrrrrr@{}}\toprule
消融 & 时间/s & 增长/\% & 95\% 区间/\% & 行程/km & 检测 & 计算/s\\\midrule
取消清后扫描 & 3081.2 & 0.15 & [-0.96, 1.32] & 11.24 & 127.4 & 0.694 \\
单方案规划 & 3131.3 & 1.78 & [0.48, 3.23] & 11.37 & 131.6 & 0.237 \\
取消测站精修 & 3136.1 & 1.93 & [0.17, 3.71] & 11.32 & 134.1 & 0.314 \\
取消概率门槛 & 3111.3 & 1.12 & [0.67, 1.66] & 11.19 & 134.2 & 0.676 
\\
\bottomrule\end{tabular}\end{table}
\begin{table}[H]\centering
\caption{消融子集的清除结果与完整时间分布}
\begin{tabular}{@{}lrrrrr@{}}\toprule
策略 & 清除比例 & 全清 & 时间均值/s & P90/s & 每源耗时/s\\\midrule
完整 joint & 1.000 & 20/20 & 3076.7 & 3424.0 & 240.3 \\
取消清后扫描 & 1.000 & 20/20 & 3081.2 & 3333.5 & 240.4 \\
单方案规划 & 1.000 & 20/20 & 3131.3 & 3444.7 & 244.3 \\
取消测站精修 & 1.000 & 20/20 & 3136.1 & 3415.2 & 245.2 \\
取消概率门槛 & 1.000 & 20/20 & 3111.3 & 3461.8 & 242.9 
\\
\bottomrule\end{tabular}\end{table}

\subsection{文件清单及复现命令}
汇总数据保存在 \texttt{B\_locator/Q3/results/}；逐动作记录和审计所需的完整输入保存在
\texttt{B\_locator/Q3/research\_archive/out/paper\_q3/}，并附有代码、期望场及场景参数的校验信息。图表从同一组数据生成，完整复现命令如下。
\begin{lstlisting}[language=bash,caption={从仓库根目录复现问题三}]
python -B B_locator/Q3/reproduce.py --audit-only
python -B B_locator/Q3/reproduce.py --out /tmp/q3-reproduction
python B_locator/Q3/research_archive/scripts/make_paper_tables.py
python B_locator/Q3/research_archive/scripts/make_paper_figures.py
python B_locator/Q3/research_archive/scripts/make_paper_trajectory_atlas.py
python B_locator/Q3/research_archive/src/p3_selftest.py --out /tmp/q3-selftest.txt
python B_locator/Q3/research_archive/src/p3_adaptive_checks.py
python B_locator/Q3/research_archive/src/p3_joint_checks.py
bash B_locator/paper/build.sh
\end{lstlisting}
实验与绘图分别使用 NumPy、Matplotlib，论文以 XeLaTeX 或 Tectonic 编译。

环境与验证说明见附件 \path{Q3/research_archive/PAPER_REPRODUCE.md}；全部 120 组五策略轨迹见 \path{Q3/figures/supplementary_trajectories.pdf}。

\clearpage
\section{问题三最终策略核心代码}\label{app:q3code}
问题三的完整可运行代码、输入数据和结果汇总均放在附件
\texttt{Q3/} 中。为保持附录紧凑，这里只列出最终
\textrm{joint} 策略的核心调度代码；其中调用的公共定位、覆盖和模拟器模块可在附件
\texttt{code/} 中按同名文件取得，文件校验值见附件的清单。
\lstinputlisting[language=python,caption={问题三最终 \textrm{joint} 策略核心代码}]{../Q3/code/p3_joint.py}
